\section{Existence of equilibrium states}

\begin{theorem}[Ground and KMS states]
\label{thm:states}
The infinite-volume system \((\A,\tau)\), and its invariant subalgebra
\((\A^G,\tau|_{\A^G})\), have at least one ground state.  For every
\(\beta\in(0,\infty)\), they have at least one \(\beta\)-KMS state.
\end{theorem}

\begin{proof}
The state space \(S(\A)\) is weak-* compact.  Choose an exhaustion
\(\Lambda_n\uparrow\Gamma\), a normalized ground vector of \(H_{\Lambda_n}\),
and extend its vector state from \(\A_{\Lambda_n}\) to \(\A\).  Gauge-average
the extension if necessary.  A subnet converges weak-* to a state
\(\omega_\infty\).  If \(A\) is local, then for all sufficiently large
\(n\), the finite-volume ground-state variational inequality gives
\[
 -i\,\omega_n(A^*\delta_{\Lambda_n}(A))
 =\omega_n\!\left(A^*[H_{\Lambda_n},A]\right)\geq0.
\tag{5.1}\label{eq:groundineq}
\]
For local \(A\), the commutator stabilizes to
\(\delta(A)=i\sum_{X:X\cap\supp A\ne\varnothing}[\Phi(X),A]\).
Passing to the weak-* limit proves the ground-state condition
\(-i\omega_\infty(A^*\delta(A))\geq0\) on the local core, hence
\(\omega_\infty\) is a ground state.

For finite \(\beta\), take the finite Gibbs state
\[
 \omega_{n,\beta}(A)=
 \frac{\operatorname{Tr}_{\Lambda_n}
 (e^{-\beta H_{\Lambda_n}}A)}
 {\operatorname{Tr}_{\Lambda_n}(e^{-\beta H_{\Lambda_n}})}
\tag{5.2}\label{eq:gibbs}
\]
and any extension to \(\A\), again gauge-averaged if desired.  Weak-*
compactness supplies a convergent subnet \(\omega_\beta\).
To pass the KMS identity without assuming complex-time volume convergence,
smear a local \(A\) with a Gaussian \(f\):
\[
 A_f=\int_{\mathbb R}f(t)\tau_t(A)\,dt,\qquad
 A_{f,n}=\int_{\mathbb R}f(t)\tau_t^{\Lambda_n}(A)\,dt.
\tag{5.3}\label{eq:smear}
\]
Theorem~\ref{thm:dynamics}, isometry, and the Gaussian tail imply
\(A_{f,n}\to A_f\) in norm; the same holds after the imaginary translation
of the entire Gaussian smear.  The finite Gibbs KMS identity
\[
 \omega_{n,\beta}(A_{f,n}B)
 =\omega_{n,\beta}\!\left(B\tau_{i\beta}^{\Lambda_n}(A_{f,n})\right)
\tag{5.4}\label{eq:finiteKMS}
\]
therefore passes to the subnet limit.  Gaussian-smeared local elements are
a norm-dense entire analytic subalgebra.  Equation
\eqref{eq:finiteKMS} on this core is equivalent to the full
\(\beta\)-KMS condition, so \(\omega_\beta\) is a KMS state.
Restriction gives states on \(\A^G\); alternatively, gauge averaging
preserves both ground and KMS conditions because the gauge action commutes
with \(\tau\).
\end{proof}

\section{Numeric twin}

The companion
\[
 \texttt{experiments/tfpt-discovery/lieb\_robinson\_witness\_probe.py}
\]
uses the exact nearest-neighbour physical-edge sector of the assembled
model while comparing it against the complete class constant
\eqref{eq:TFPTvelocity}.  This is the same finite cone protocol as the
earlier IR-universality witness.  It reports:
\[
\begin{array}{c|c}
\toprule
\text{quantity}&\text{value}\\
\midrule
v_{\rm LR}\text{ (proved class bound)}&387.639069990831\\
v_{\rm measured}\text{ (commutator threshold }0.02)&2.000000000000\\
\text{outside-cone leakage}&8.238\times10^{-3}\\
\bottomrule
\end{array}
\tag{6.1}\label{eq:numericvelocity}
\]
For the one-site projector, \(t=10^{-3}\), and nested open volumes, the
specialization of \eqref{eq:convergence} gives
\[
\begin{array}{c|cc}
\toprule
b&\norm{\tau_t^{\Lambda_b}(A)-\tau_t^{\Lambda_{64}}(A)}
&S_2(b-2)(e^{2e\kappa t}-1)\\
\midrule
8&0&9.246214927512\times10^{-1}\\
16&0&3.279645827421\times10^{-2}\\
32&0&2.164381663985\times10^{-5}\\
\bottomrule
\end{array}
\tag{6.2}\label{eq:numericconvergence}
\]
The zeros are numerical precision statements, not analytic finite-speed
claims.  A \(1/r\) all-to-all mutant, which is outside
\eqref{eq:LRhyp}, gives \(1.048309991863\times10^{-4}\) at \(b=32\),
violating the inapplicable finite-range line by a factor \(4.843462\).
The Gauss-positive control has
\(\norm{[A,P_G]}=\norm{P_GAP_G-AP_G}=0\); the constraint-violating
embedding has \(\norm{[A_{\rm bad},P_G]}=1\).  All eight checks pass.

\section{Combined theorem and exact boundary}

\begin{theorem}[Mandatory thermodynamic-dynamics leg]
\label{thm:combined}
For the TFPT master-Hamiltonian class \eqref{eq:H} with the frozen
finite-dimensional local factors and coefficients of
Table~\ref{tab:terms}:
\begin{enumerate}[label=\textup{(\roman*)}]
\item the finite-volume dynamics satisfy the uniform Lieb--Robinson bound
\eqref{eq:LR} with the explicit estimate
\(v_{\rm LR}=32086e/225\);
\item the thermodynamic-limit dynamics exist in norm on all local
observables and define a strongly continuous automorphism group;
\item the group restricts consistently to the Gauss-invariant quasilocal
algebra, with projector compatibility in the precise form
\eqref{eq:projectorcompat}; and
\item at least one thermodynamic ground state and at least one
\(\beta\)-KMS state for every \(\beta>0\) exist.
\end{enumerate}
No continuum limit is used in any part of the proof.
\end{theorem}

\begin{proof}
Part (i) is Theorem~\ref{thm:LRclassical} together with
\eqref{eq:constants}.  Part (ii) is Theorem~\ref{thm:dynamics}, including
the explicit norm-Cauchy estimate \eqref{eq:convergence}.  Part (iii) is
Proposition~\ref{prop:gauss}; the warning following it excludes the false
physical-Hilbert-space embedding.  Part (iv) is Theorem~\ref{thm:states}.
\end{proof}

\begin{center}
\setlength{\fboxsep}{10pt}
\fcolorbox{red!65!black}{red!3}{
\begin{minipage}{0.91\textwidth}
\textbf{BOUNDARY BOX --- what this does not give.}

\smallskip
This theorem does \emph{not} give uniqueness of the ground or KMS state;
phase structure and spontaneous symmetry breaking remain undecided.  It
does \emph{not} prove that the finite-volume spectral gap survives the
thermodynamic limit.  It does \emph{not} establish an interacting
continuum/scaling limit, Lorentz invariance, a common infrared light cone,
or IR universality.  Those are precisely the remaining T3/T5 contents.
The result proved here is the mandatory dynamics leg only: a uniformly
local Hamiltonian class, a volume-independent Lieb--Robinson estimate,
norm thermodynamic-limit dynamics, a compatible gauge-invariant
quasilocal algebra, and existence (not uniqueness) of equilibrium states.
\end{minipage}}
\end{center}

\section*{Conclusion}
\addcontentsline{toc}{section}{Conclusion}

The amended contract makes the continuum limit optional but the
thermodynamic limit mandatory.  The latter is now proved at the
Hamiltonian-class level for the concrete finite-group, finite-local-space
TFPT assembly.  The proof is insensitive to finite-volume boundary shape
and does not infer thermodynamics from a finite scan.  Its explicit
velocity is conservative, while the numerical witness checks the expected
cone, the norm-convergence inequality, and both destructive mutants.
The boundary box prevents this dynamics theorem from being misread as a
gap, phase-uniqueness, or IR-universality theorem.

\begin{thebibliography}{9}

\bibitem{LR1972}
E.~H.~Lieb and D.~W.~Robinson,
\emph{The finite group velocity of quantum spin systems},
Commun.\ Math.\ Phys.\ \textbf{28} (1972), 251--257,
\href{https://doi.org/10.1007/BF01645779}
{doi:10.1007/BF01645779}.

\bibitem{NS2006}
B.~Nachtergaele and R.~Sims,
\emph{Lieb--Robinson bounds and the exponential clustering theorem},
Commun.\ Math.\ Phys.\ \textbf{265} (2006), 119--130,
\href{https://doi.org/10.1007/s00220-006-1556-1}
{doi:10.1007/s00220-006-1556-1}.

\bibitem{NS2010}
B.~Nachtergaele and R.~Sims,
\emph{Lieb--Robinson bounds in quantum many-body physics},
Contemp.\ Math.\ \textbf{529} (2010), 141--176,
\href{https://doi.org/10.1090/conm/529/10429}
{doi:10.1090/conm/529/10429}.

\bibitem{BR}
O.~Bratteli and D.~W.~Robinson,
\emph{Operator Algebras and Quantum Statistical Mechanics II},
second edition, Springer, Berlin (1997), Theorem~6.2.11.

\bibitem{GrundlingLledo}
H.~Grundling and F.~Lled\'o,
\emph{Local quantum constraints},
Rev.\ Math.\ Phys.\ \textbf{12} (2000), 1159--1218.

\bibitem{GrundlingRudolph}
H.~Grundling and G.~Rudolph,
\emph{QCD on an infinite lattice},
Commun.\ Math.\ Phys.\ \textbf{349} (2017), 1163--1202,
\href{https://doi.org/10.1007/s00220-013-1674-5}
{doi:10.1007/s00220-013-1674-5}.

\end{thebibliography}

\end{document}
